\section{Spectral capacity, energy memory, and thermal transport}
\label{v:sec:capacidad-espectral-termica}

The positive map between temperature and decision energy is constructed here
on the graded carrier and its marked reference. This section retains the
energy origin invisible to
normalisation, and composes that reading with the calendar, area, gravitation,
and the thermal distributions of this article. The causal chain remains
\[
 \mathrm{APP}\longrightarrow\mathrm{TRIT}\longrightarrow\mathrm{TPK}
 \longrightarrow X_{\rm enr}
 \longrightarrow\mathcal C_{\rm cont}^{\rm disc}
 \longrightarrow\text{thermal carrier}
 \longrightarrow\text{physical recognition}.
\]
The mathematical and physical constants used here are HMT--MD outputs or
later representations of that state. No metrological number selects seeds,
sectors, or spectrum. In particular, the scalar
$1.380649\times10^{-23}$ is first obtained as a marked dimensionless
intensity; its composition with the energy and thermal charts is proved only
afterwards.

\subsection{Graded carrier and two preserved coordinates}

Let $\mathscr H_{N,L}$ be five finite spaces, with
\begin{equation}
 \dim\mathscr H_{0,0}=1,\quad
 \dim\mathscr H_{1,2}=380,\quad
 \dim\mathscr H_{2,0}=1,\quad
 \dim\mathscr H_{2,1}=24,\quad
 \dim\mathscr H_{2,2}=624.
 \label{v:eq:fibras-termicas-NL}
\end{equation}
Their direct sum $\mathscr H_{\rm th}$ simultaneously retains the prefix
degree $N$ and the incidence-memory degree $L$. Its bivariate Hilbert
polynomial is
\begin{equation}
 Z(u,v)=1+380uv^2+u^2(1+24v+624v^2).
 \label{v:eq:polinomio-graduado-termico}
\end{equation}
Thus,
\[
 Z(u,1)=1+380u+649u^2,\qquad
 Z(1,v)=2+24v+1004v^2.
\]
The equality $649=1+24+624$ describes the $L$-fibres in degree $N=2$;
the independent identity $649=1+2\cdot324$ describes the involutive
decomposition of the carrier. These meanings are not interchanged.

The dimensions $380$ and $649$ arise from distinct operations. The first is
the off-diagonal part of the endomorphism on the capacity-$20$ carrier; the
second joins the scalar direction to two orientations of the $324$-route
sector. To make the latter construction explicit, let
\[
 X_\times=(\mathbb Z/9\mathbb Z)^2\times\{N,E,S,O\},
 \qquad |X_\times|=324,
\]
and let $\mathfrak e:X_\times\to\mathcal S_\times$ be the six-window
emitter, with $26$ signatures. Its fibres have size $8$ for $24$ signatures,
$24$ for $555555$, and $108$ for $000000$. With
$I_{24}=\{1,\ldots,24\}$, construct the microscopic tree
\begin{equation}
 \widetilde{\mathcal T}=\{\Omega\}\sqcup\{a_i:i\in I_{24}\}
 \sqcup\{b_{i,y}:i\in I_{24},\ y\in X_\times\},
 \qquad |\widetilde{\mathcal T}|=7801.
 \label{v:eq:arbol-microscopico-7801}
\end{equation}
The reading $p(b_{i,y})=b_{i,\mathfrak e(y)}$, fixing $\Omega$ and the
$a_i$, has image $\mathcal T$ of size
$1+24+24\cdot26=649$. If
$m_s=|\mathfrak e^{-1}(s)|$, the measure
\begin{equation}
 \nu_{\rm mic}(\Omega)=\nu_{\rm mic}(a_i)=1,
 \qquad \nu_{\rm mic}(b_{i,y})=m_{\mathfrak e(y)}^{-1}
 \label{v:eq:medida-arbol-microscopico}
\end{equation}
gives every fibre of $p$ mass one. Rates are $r>0$ in both directions
between $\Omega$ and $a_i$, $r/m_{\mathfrak e(y)}$ from $a_i$ to
$b_{i,y}$, and $r$ in the reverse direction.

On $\ell^2(\mathcal T)$ with counting measure and
$\ell^2(\widetilde{\mathcal T},\nu_{\rm mic})$, the lift
\begin{equation}
 J_p g=g\circ p
 \label{v:eq:levantamiento-fibras-Jp}
\end{equation}
is an isometry from a class or signature to the weighted microscopic tree.
Indeed,
\[
 \|J_pg\|_{\nu_{\rm mic}}^2
 =\sum_{z\in\mathcal T}|g(z)|^2
   \sum_{x:p(x)=z}\nu_{\rm mic}(x)=\|g\|^2,
 \qquad J_p^*J_p=I.
\]
Its adjoint is the weighted average
\begin{equation}
 (J_p^*f)(z)=\sum_{x:p(x)=z}\nu_{\rm mic}(x)f(x),
 \label{v:eq:adjunto-fibras-Jp}
\end{equation}
not the lift $J_p$: their domains are reversed, and one copies a class
function whereas the other averages its fibre. Detailed balance
$\nu_{\rm mic}(a_i)r/m_{\mathfrak e(y)}
=\nu_{\rm mic}(b_{i,y})r$ and the rate sums prove, for the microscopic
and aggregated generators,
$\widetilde L J_p=J_pL_{\mathcal T}$. Balance within a fibre does not by
itself determine weights between grades; the boundary currents constructed
below do so.

The variable
\begin{equation}
 Q:=N-L
 \label{v:eq:carga-memoria-Q}
\end{equation}
separates states that one marginal may identify. For example, knowing
$Z(u,1)$ does not recover the distribution of $L$. The five cells in
\eqref{v:eq:fibras-termicas-NL} and their multiplicities reconstruct the
specified joint carrier; they do not reconstruct coherences discarded by a
compression. This is the exact domain of the reading.

\subsection{Positive lift of energy memory}

On the five cells, $Q$ takes the values $0,-1,2,1,0$. Introduce a memory
register $R$ with levels $0,1,2,3$ and the isometry
\begin{equation}
 V_{\rm mem}:\mathscr H_{N,L}\longrightarrow
 \mathscr H_{N,L}\otimes\mathbb C^4,\qquad
 V_{\rm mem}\xi=\xi\otimes e_{N-L+1}.
 \label{v:eq:isometria-memoria-positiva}
\end{equation}
The unit shift turns the signed charge into a nonnegative level without
losing it. If $R e_j=j e_j$, then
\begin{equation}
 V_{\rm mem}^*(L\otimes I+I\otimes R)V_{\rm mem}=N+I.
 \label{v:eq:balance-positivo-memoria}
\end{equation}
On each cell this is simply $L+(N-L+1)=N+1$, but the operator form retains
degeneracies, superpositions within each fibre, and every purification on an
auxiliary system.

For an energy scale $\epsilon_0>0$,
\begin{equation}
 V_{\rm mem}^*\epsilon_0(L\otimes I+I\otimes R)V_{\rm mem}
 =\epsilon_0(N+I)
 \label{v:eq:energia-lift-memoria}
\end{equation}
proves exact energy conservation in the lift. A filter retaining one branch
must also retain its success marker; otherwise the conditioned density hides
the preparation cost. The lift is reversible on its image and has zero
logical cost; discarding $R$ is a distinct operation governed by the Landauer
balance of the preceding section.

\subsection{Decimal spectrum produced by the memory loop}

The $8{:}1$ incidence decomposition determines the coupling
\[
 \mathcal Q=\frac13
 \begin{pmatrix}\sqrt8 I&-I\\ I&\sqrt8 I\end{pmatrix}.
\]
The oriented return loop is
\[
 H_1=\begin{pmatrix}2&1\\1&1\end{pmatrix},\qquad
 H_1H_1^T=\begin{pmatrix}5&3\\3&2\end{pmatrix}.
\]
In the declared positive Gaussian preparation, the normalised visible
covariance is
\begin{equation}
 W=\frac{8I+H_1H_1^T}{9},\qquad
 \det W=\frac{121}{81},\qquad
 \nu_W=\sqrt{\det W}=\frac{11}{9}.
 \label{v:eq:williamson-decimal}
\end{equation}
The one-mode spectral theorem therefore gives
\begin{equation}
 r=\frac{\nu_W-1}{\nu_W+1}=\frac1{10},\qquad
 \rho_r=(1-r)\sum_{j\ge0}r^j\,|j\rangle\langle j|.
 \label{v:eq:estado-geometrico-decimal}
\end{equation}
The decimal ratio is obtained from the loop covariance; it is not entered as
a target numerical base. The symbol $\nu_W$ is the Williamson symplectic
eigenvalue and is distinct from the microscopic measure $\nu_{\rm mic}$ in
\eqref{v:eq:medida-arbol-microscopico}. The control $H_0=I$ gives $r=0$,
so the spectrum depends effectively on the selected holonomy.

Euclidean division $j=3n+b$, $0\le b<3$, defines the unitary
\begin{equation}
 U_3|j\rangle=|n\rangle\otimes|b\rangle,\qquad
 U_3\rho_rU_3^*=\rho_{r^3}\otimes
 \frac{\operatorname{diag}(1,r,r^2)}{1+r+r^2}.
 \label{v:eq:agrupamiento-tritico-espectral}
\end{equation}
The proof is termwise:
$(1-r)r^{3n+b}=(1-r^3)(r^3)^n r^b/(1+r+r^2)$. Hence,
\begin{equation}
 q=r^3=10^{-3},\qquad
 \sigma_r=\frac{\operatorname{diag}(100,10,1)}{111}.
 \label{v:eq:razon-milenaria-y-residuo}
\end{equation}
The residue retains the three intermediate positions. If
$S|j\rangle=|j+1\rangle$, its transport is
\[
 U_3SU_3^*=I\otimes(|1\rangle\langle0|+|2\rangle\langle1|)
 +S_n\otimes|0\rangle\langle2|,
\]
so the third advance returns the residue and increments the quotient. Phase
return and state return are not conflated.

With $H_W=\epsilon_a(N+\tfrac12)$,
\begin{equation}
 U_3H_WU_3^*
 =\epsilon_a\left(3N_n\otimes I+I\otimes R+\tfrac12I\right).
 \label{v:eq:hamiltoniano-agrupado-decimal}
\end{equation}
The quotient factor has spacing $3\epsilon_a$, while temperature is fixed:
\[
 e^{-\beta_W\epsilon_a}=10^{-1},\qquad
 e^{-\beta_W3\epsilon_a}=10^{-3},\qquad
 k_BT_W=\frac{\epsilon_a}{\log10}.
\]

\subsection{Marked section and complete evaluation}

The renewal in~\eqref{v:eq:tupla-renovacion-capacidad} supplies origin $23$,
and threefold refinement supplies $23,26,29$. With
$\lambda_j/\lambda_0=r^j$ and the carrier multiplicities,
\begin{equation}
 \boxed{
 \mathcal B
 =\frac{\lambda_{23}+380\lambda_{26}+649\lambda_{29}}{\lambda_0}
 =10^{-23}\left(1+\frac{380}{1000}+
                    \frac{649}{1000^2}\right)
 =1.380649\times10^{-23}.}
 \label{v:eq:evaluacion-termica-completa}
\end{equation}
The residue $b=2$ alone would select $2,5,8,\ldots$ and does not determine
threshold $23$; the latter comes from renewal time and the capacity index.
The mantissa is retained in the normalised density through the sector marker;
the selection weight additionally retains $10^{-23}$. Thus neither the
mantissa nor the exponent is taken from an external table.

Let
$\mathcal D=\bigoplus_{n=0}^2\mathcal D_n$, with
$(d_0,d_1,d_2)=(1,380,649)$,
$\dim\mathcal D_n=d_n$, and $N|_{\mathcal D_n}=nI$. Define the capacity
operator and spectral reference by
\begin{equation}
 C=23I+3N,\qquad
 \eta=\bigoplus_{n=0}^2r^{23+3n}I_{\mathcal D_n},\qquad
 b_*:=\operatorname{Tr}\eta=\mathcal B.
 \label{v:eq:referencia-termica-eta}
\end{equation}
Reference $\eta$ retains grades $23,26,29$ separately and remains fixed
when a trial state varies. For $x=e^{-\epsilon_0/\tau}$, the relative
partition function and state are
\begin{equation}
 Z_N(x)=1+380x+649x^2,\qquad
 \rho_N(x)=\frac{x^N}{Z_N(x)},\qquad
 p_0=\frac1{Z_N(x)}.
 \label{v:eq:particion-portador-termico}
\end{equation}
The marker $p_0$ recovers $Z_N$ after normalisation. At $q=10^{-3}$,
$Z_N(q)=1.380649$, $\eta=10^{-23}q^N$, and
$\eta/b_*=\rho_N(q)$; the marked origin produces
\eqref{v:eq:evaluacion-termica-completa}. The free energy and mean degree are
\begin{equation}
 F_N(\tau)=-\tau\log Z_N(x),\qquad
 \langle N\rangle_x=x\frac{d}{dx}\log Z_N(x)
 =\frac{380x+1298x^2}{Z_N(x)}.
 \label{v:eq:energia-libre-grado-medio}
\end{equation}

\subsection{Boundary currents and rigidity of the sector origin}

The areal incidence supplies $18+18$ tritic links. Let
$u=N_{90}$, $v=N_{120}$, $a=u+v$, and
\[
 F=3u+4v,\qquad H_R=\epsilon_0F,
 \quad \epsilon_0=30\Delta\tau_R>0,
 \quad \tau_R=\frac{\hbar c}{2\pi R}.
\]
Provenance remains in $(u,v)$: the same area may correspond to different
energies. On one link,
$T^+=|1\rangle\langle0|+|2\rangle\langle1|$ and $T^-=(T^+)^*$ satisfy
$[n,T^\pm]=\pm T^\pm$. For $e,f=1,\ldots,18$, define
\[
 J_{ef}=T^+_{120,f}T^-_{90,e},\qquad
 B_c=\left(\sum_eT^+_{90,e}\right)^3
     \left(\sum_fT^-_{120,f}\right)^2.
\]
The rule $[F,AB]=[F,A]B+A[F,B]$ proves
\[
 [F,J_{ef}]=J_{ef},\quad[a,J_{ef}]=0,
 \qquad [F,B_c]=B_c,\quad[a,B_c]=B_c.
\]
The occupation sequence
$(0,2)\to(3,0)\to(2,1)$ has primitive energies $8,9,10$ and areas
$2,3,3$: the same energy increment retains two distinguishable areal
actions.

The coefficients $c_j=[z^j](1+z+z^2)^{18}$ count occupations in each
channel: $c_0=1,c_1=18,c_2=171,c_3=1122$. The boundary character
\[
 P(x)=(1+x^3+x^6)^{18}(1+x^4+x^8)^{18}
\]
gives $g_7=324,g_8=171,g_9=1122,g_{10}=3078$. For ranks
$(1,380,649)$, the first admissible triple of consecutive energies starts
at $8$. Earlier starts are excluded because, using
$g_0=1$, $g_1=g_2=g_5=0$, $g_3=g_4=18$, $g_6=171$, $g_7=324$, and
$g_8=171$, every earlier triple fails at least one required rank.

The connected compression is constructed by selecting one $(0,2)$ state,
$380$ states $(3,0)$, and $649$ distinct neighbours under the $J_{ef}$.
Every occupation of sum three has at least two predecessors of sum two,
and every predecessor has at most $18$ successors. Hence the $380$ states
produce at least $\lceil760/18\rceil=43$ distinct predecessors and
$18\cdot43=774$ neighbours $(2,1)$, from which $649$ may be selected.
The current $B_c$ connects the initial state to the $380$ intermediate
states, and every final neighbour retains a selected predecessor; the
resulting graph is connected. Thus the compressed subspace, its labels,
and its currents are constructed rather than merely postulated.

The boundary compression realises degrees $N=0,1,2$ at energies
$8\epsilon_0,9\epsilon_0,10\epsilon_0$. Let
$\mathcal D_n=\ker(N-nI)$,
$P_F=\mathbf1_{\{N=1\}}$,
$B_N:\mathcal D_0\to\mathcal D_1$, and
$J_N:\mathcal D_1\to\mathcal D_2$ be the nonzero compressed currents;
their adjoints traverse the edges downwards. Then
\begin{equation}
 [N,B_N]=B_N,\quad [N,J_N]=J_N,\qquad
 [P_F,B_N]=B_N,\quad [P_F,J_N]=-J_N.
 \label{v:eq:conmutadores-corrientes-termicas}
\end{equation}
For the affine diagonal correction
\[
 H_\delta=\epsilon_0(8I+N)+\delta P_F
\]
the commutators give
\[
 [H_\delta,B_N]=(\epsilon_0+\delta)B_N,
 \qquad
 [H_\delta,J_N]=(\epsilon_0-\delta)J_N.
\]
Requiring both currents to retain the boundary jump $\epsilon_0$, or merely
requiring both jumps to be equal, forces
\begin{equation}
 \delta=0.
 \label{v:eq:rigidez-desfase-corrientes}
\end{equation}
The conclusion classifies this sector offset; it does not claim a
classification of the entire commutant of the Hamiltonian.

Assigning rate $xr_e$ to every ascending edge and rate $r_e>0$ to its
descending mate, with the same $r_e$ in both directions, gives detailed
balance with respect to weights $x^N$. Connectivity makes the stationary
distribution unique on the compressed carrier and recovers
\eqref{v:eq:particion-portador-termico}. This dynamics is a Markov preparation
semigroup, not isolated unitary evolution.

\subsection{Thermal equivalence with origin and marker}

Let $(\mathscr H_i,H_i,\tau_i,e_i,P_i)$ be two realisations, with
$\tau_i>0$, origin $e_i$, support projector $P_i$, and an isometric
isomorphism $J:P_1\mathscr H_1\to P_2\mathscr H_2$. The relevant thermal
equivalence is
\begin{equation}
 J\frac{H_1-e_1I}{\tau_1}
 =\frac{H_2-e_2I}{\tau_2}J,\qquad JP_1=P_2J.
 \label{v:eq:equivalencia-termica-centrada}
\end{equation}
It requires neither $H_1=H_2$ nor $\tau_1=\tau_2$. It implies
$Jf((H_1-e_1I)/\tau_1)=f((H_2-e_2I)/\tau_2)J$ for every bounded spectral
function $f$, and transports the centred Gibbs states. For observables that
depend on absolute energy, the $e_i$ and scale must also be transported:
normalisation does not license their removal.

In the preceding compression there is a based isometry $\iota_N$ such that
\[
 H_R\iota_N=\iota_N\epsilon_0(8I+N).
\]
The current $J_N:\mathcal D_1\to\mathcal D_2$ and the isometry
$\iota_N$ are distinct operators. For the chronological reading, set
\begin{equation}
 C=23I+3N,\qquad
 H_{\rm cap}=\frac{\epsilon_a\log10}{\log3}C,
 \qquad \tau_{\rm clk}=\frac{\epsilon_a}{\log3}.
 \label{v:eq:lector-capacidad-escalas}
\end{equation}
The relative Williamson section instead uses
$H_{W,{\rm rel}}=\epsilon_a C$ and
$\tau_W=\epsilon_a/\log10$. Hence
$H_{\rm cap}/\tau_{\rm clk}=H_{W,{\rm rel}}/\tau_W=C\log10$,
although their energy and temperature scales differ.

With $\tau_*=\epsilon_0/\log1000$, the centred comparison is
\begin{equation}
 \frac{H_R-8\epsilon_0I}{\tau_*}\iota_N
 =\iota_N\frac{H_{\rm cap}-23\tau_{\rm clk}\log10\,I}
                    {\tau_{\rm clk}}
 =\iota_N(\log1000)N.
 \label{v:eq:transporte-termico-centrado}
\end{equation}
The equality retains degrees, energies, marker, and origin without
identifying $\epsilon_0$ with $\epsilon_a$: the former belongs to the
boundary compression and the latter to the chronological section. The
$1/10$ factor between uncentred traces and the $\log10$ shift are section
data, not defects of the density.

\subsection{Area, provenance, and thermal cost}

In the two areal channels, let $N_{90},N_{120}\ge0$ and
\begin{equation}
 a=N_{90}+N_{120},\qquad D=90N_{90}+120N_{120}.
 \label{v:eq:ocupacion-procedencia-90-120}
\end{equation}
They are recovered exactly as
\begin{equation}
 N_{90}=4a-\frac D{30},\qquad
 N_{120}=\frac D{30}-3a.
 \label{v:eq:inversion-procedencia-90-120}
\end{equation}
Area may read $a$, while the sector Hamiltonian retains provenance:
\begin{equation}
 H_R=\tau_R\Delta D
 =\tau_R\Delta(90N_{90}+120N_{120}).
 \label{v:eq:hamiltoniano-area-procedencia}
\end{equation}
Two states with equal $a$ and different channel distributions have the same
total-occupation reading and different cost. Reduction to one channel, its
purification, and the finite free-energy identity apply to the reduced state
already constructed; relative entropy measures the deficit from the reference
Gibbs state. This composition explains why origin and provenance remain
necessary when relating area, entanglement, and the Barbero--Immirzi
functional.

\subsection{Composition with action, gravitation, and the electron}

The length $L_\alpha$ and oriented action $\hbar_a$ fix
\begin{equation}
 \epsilon_a=\frac{\hbar_ac}{L_\alpha},\qquad
 G_a=\frac{c^3L_\alpha^2}{\hbar_a},\qquad
 k_B\Theta_{{\rm clk},a}\log3=\epsilon_a.
 \label{v:eq:triple-accion-gravedad-temperatura}
\end{equation}
Multiplying the three identities eliminates the action section and gives
\begin{equation}
 \boxed{G_a k_B\Theta_{{\rm clk},a}\log3=c^4L_\alpha.}
 \label{v:eq:identidad-gravedad-temperatura}
\end{equation}
The identity is obtained in one orientation $a$; $G_a$ is not taken from
metrology. The symbol $L_\alpha$ denotes length, not the degree $L$ of
\eqref{v:eq:fibras-termicas-NL}.

If $R_{\rm act}=\hbar_{\rm pre}/\hbar_{\rm ret}$ is the action ratio of the
electron return, the same composition gives
\begin{equation}
 R_{\rm act}
 =\frac{\Theta_{{\rm clk,pre}}}{\Theta_{{\rm clk,ret}}}
 =\frac{G_{\rm ret}}{G_{\rm pre}}.
 \label{v:eq:retorno-accion-termico-gravitatorio}
\end{equation}
This equality transports the thermal section; it neither uses $k_B$ to
select the route nor rederives the mass law in this article. For an already
individualised electron energy $E_e$, $y_e=E_e/\epsilon_{\rm ret}$ satisfies
\[
 \beta_{\rm clk}E_e=(\log3)y_e,\qquad
 \frac{G_{\rm ret}m_e^2}{\hbar_{\rm ret}c}=y_e^2.
\]
The energy origin affects mass and radius even when it cancels from
normalised probabilities.

\subsection{Marked thermometric bridge and effective identity}

Reference $\eta$ in~\eqref{v:eq:referencia-termica-eta} must remain
separate from the state on which energy is measured. Let $\mathcal H$ be a
finite-dimensional space, let $\rho$ be a density matrix on $\mathcal H$,
and let $H=H^*$ be a Hamiltonian with registered origin. On
$\widehat{\mathcal D}=\mathcal D\oplus\mathbb C f$, define
\begin{equation}
 \widehat\eta=\eta\oplus(1-b_*)|f\rangle\langle f|,
 \qquad \Omega_\rho=\widehat\eta\otimes\rho,
 \label{v:eq:referencia-marcada-normalizada}
\end{equation}
and let $P$ project onto the summand $\mathcal D$, extended by the identity
on $\mathcal H$. With
$s(\rho)=-\operatorname{Tr}(\rho\log\rho)$ and logarithms restricted to
the support, the two local readings are
\begin{align}
 \mathcal S_P(\rho)
 &=-\operatorname{Tr}\!\left[P\Omega_\rho
                 (I\otimes\log\rho)\right]
   =b_*s(\rho),\label{v:eq:lector-entropico-marcado}\\
 \mathcal E_P(\rho)
 &=\frac{\operatorname{Tr}\!\left[P\Omega_\rho(I\otimes H)\right]}
         {\operatorname{Tr}(P\Omega_\rho)}
   =\operatorname{Tr}(\rho H).
 \label{v:eq:lector-energetico-condicionado}
\end{align}
Indeed, $P\Omega_\rho=\eta\otimes\rho$ and
$\operatorname{Tr}(P\Omega_\rho)=b_*$. Trace factorisation gives
$b_*s(\rho)$ and $b_*\operatorname{Tr}(\rho H)$ in the two numerators;
dividing the latter by $b_*>0$ proves the energy reading. Thus the fixed
reference, information relative to its marker, and energy per conditioned
state are all retained.

Work in transported positive bases of the energy, temperature, and entropy
lines, with $u_{\rm ent}=u_E/u_\Theta$. In coordinate formulas, $H$ expresses
energy in $u_E$ and $\mathcal S_P$ expresses entropy in $u_{\rm ent}$. The
realisation rule comprises four operations: receive the fixed
reference~\eqref{v:eq:referencia-marcada-normalizada};
evaluate~\eqref{v:eq:lector-entropico-marcado};
evaluate~\eqref{v:eq:lector-energetico-condicionado}; and define temperature
as conjugate to these two functionals. The bases are transported with the
operators, without a target temperature input.

\begin{theorem}[Thermometric derivative of the marked reading]
\label{v:thm:derivada-termometrica-marcada}
If $H$ is nonscalar and $\rho_\beta=e^{-\beta H}/Z(\beta)$ for $\beta>0$,
then
\begin{equation}
 \frac{d\mathcal S_P}{d\mathcal E_P}=b_*\beta,
 \qquad \Theta_P=\frac1{b_*\beta},
 \qquad B_*(\Theta_Pu_\Theta)=\beta^{-1}u_E,
 \label{v:eq:derivada-termometrica-marcada}
\end{equation}
where $B_*(u_\Theta)=b_*u_E$ is the unique positive linear map with that image
of the basis.
\end{theorem}
\begin{proof}
For $E(\beta)=\operatorname{Tr}(\rho_\beta H)$, the finite spectral sum
gives $E'=-\operatorname{Var}_{\rho_\beta}(H)<0$ and
$(\log Z)'=-E$. From
$s(\rho_\beta)=\beta E+\log Z$ one obtains $s'=\beta E'$. Since $\eta$
remains fixed, $\mathcal S_P'=b_*s'$ and $\mathcal E_P'=E'$. Their ratio
and its reciprocal prove the first two identities; the image of a positive
basis fixes a unique linear map and proves the third.
\end{proof}

The factor depends on the pair of readings. If energy and entropy are both
marked, $(b_*E,b_*s)$, or both conditioned, $(E,s)$, the derivative is
$\beta$; for the marked--conditioned pair
\eqref{v:eq:lector-entropico-marcado}--
\eqref{v:eq:lector-energetico-condicionado} it is $b_*\beta$.

The same rule has a variational characterisation. For $\Theta>0$,
\begin{equation}
 \Phi_\Theta(\rho)=\mathcal E_P(\rho)-\Theta\mathcal S_P(\rho)
 =\operatorname{Tr}(\rho H)-b_*\Theta s(\rho)
 \label{v:eq:funcional-variacional-marcado}
\end{equation}
has the unique minimiser
\begin{equation}
 \rho_\Theta^*=\frac{e^{-H/(b_*\Theta)}}
                    {\operatorname{Tr}e^{-H/(b_*\Theta)}}.
 \label{v:eq:minimo-variacional-marcado}
\end{equation}
Indeed,
\[
 \Phi_\Theta(\rho)-\Phi_\Theta(\rho_\Theta^*)
 =b_*\Theta D_{\rm rel}(\rho\Vert\rho_\Theta^*)\ge0,
\]
where
$D_{\rm rel}(\rho\Vert\sigma)
=\operatorname{Tr}[\rho(\log\rho-\log\sigma)]$; equality in
relative-entropy positivity occurs exactly when
$\rho=\rho_\Theta^*$. This proves existence and uniqueness without entering
a target temperature.

Apply the result to the capacity clock. With orientation and $\epsilon_a$
fixed, set $H_{{\rm clk},a}=H_{\rm cap}$ and
$\beta_{{\rm clk},a}=\log3/\epsilon_a$. This Hamiltonian is positive,
self-adjoint, and nonscalar, since its three distinct eigenvalues are
proportional to $23,26,29$. Then
\begin{align}
 \beta_{{\rm clk},a}H_{{\rm clk},a}
   &=C\log10,\nonumber\\
 e^{-\beta_{{\rm clk},a}H_{{\rm clk},a}}
   &=10^{-23}q^N=\eta,\qquad q=10^{-3},\nonumber\\
 \rho_{{\rm clk},a}
   &=\frac{q^N}{Z_N(q)}=\frac{\eta}{b_*}.
 \label{v:eq:estado-reloj-capacidad}
\end{align}
Reference $\eta$, independent of $\epsilon_a$, stays fixed in the first
factor, while $\rho_\beta$ varies on a second copy of $\mathcal D$ with
$H_{{\rm clk},a}$ fixed. Evaluation only afterwards at
$\beta=\beta_{{\rm clk},a}$ gives
\begin{equation}
 \Theta_{{\rm clk},P,a}=\frac{\epsilon_a}{b_*\log3},
 \qquad
 \boxed{B_*(\Theta_{{\rm clk},P,a}u_\Theta)
 =\frac{\epsilon_a}{\log3}u_E
 =\frac{h_a}{108t_0\log3}u_E.}
 \label{v:eq:identidad-transductor-numeral}
\end{equation}
The last equality uses the already fixed action coordinate
$\epsilon_a=h_a/(108t_0)$; it does not use the value of $b_*$ as a fit.
Nor does it identify $\epsilon_a$ with the boundary scale $\epsilon_0$.

If $B_{\rm clk}$ is the chronological transducer and $J_E,J_\Theta$
transport energy and the two thermometric sections of the same reference
state, respectively, uniqueness between lines gives the typed square
\begin{equation}
 B_*J_\Theta=J_EB_{\rm clk}.
 \label{v:eq:cuadrado-transductor-marcado}
\end{equation}
The two composites agree on the positive generating section by
\eqref{v:eq:identidad-transductor-numeral}; linearity extends the equality
to the entire line. This is the bridge between the spectral evaluation
in~\eqref{v:eq:evaluacion-termica-completa} and the dimensional transducer:
domains, reference, readers, and positive section determine the map before
its numerical evaluation.

\clearpage
\subsection{Consequences for occupation, radiation, and information}

The finite carrier supplies Gibbs theory with $Z_N(x)$, $\rho_N(x)$, marker
$p_0$, and the free energy~\eqref{v:eq:energia-libre-grado-medio}. The
symmetric and exterior completions developed below act on complete modes and
produce the Bose--Einstein and Fermi--Dirac laws; the $1/380/649$ carrier does
not replace them. It supplies an energy section and a preparation with a
recorded origin. In the dilute regime, the same coordinate
$x=e^{-\beta(E-\mu)}$ enters the Maxwell--Boltzmann limit already proved.

In the radiative realisation, composing $\beta=(k_B\Theta)^{-1}$ with
$k_B\Theta_{{\rm clk},a}\log3=\epsilon_a$ substitutes the calendar scale
in Planck's density without changing the count of two polarisations or their
dispersion. Its moments retain the Stefan--Boltzmann and Wien deductions; the
new section determines the source of the thermal scale and the marker that
must be retained when conditioning a branch. For area and entanglement,
\eqref{v:eq:inversion-procedencia-90-120} prevents total occupation from
erasing provenance cost. Thus occupation, radiation, free energy, area, and
gravitation receive the same thermal information without conflating their
observables.
